\documentclass[11pt,a4paper]{article}
\usepackage[margin=2.5cm]{geometry}
\usepackage{amsmath}
\usepackage{booktabs}
\usepackage[hidelinks]{hyperref}
\usepackage{parskip}

\title{NucAE: reconstructing nucleosome coverage from sparse cfDNA sequencing}
\author{}
\date{\today}

\begin{document}
\maketitle

\section{What the system does}

Cell-free DNA fragments in blood are protected from degradation by nucleosomes,
so the midpoint of a fragment marks approximately where a nucleosome sat. Counting
those midpoints along the genome gives a coverage track whose peaks are nucleosome
positions.

At high sequencing depth this track is clean. At the depth a clinical assay can
afford it is mostly noise. The model takes the noisy track and reconstructs what
the deep one would have looked like.

\begin{center}
\begin{tabular}{lll}
\toprule
& \textbf{input} & \textbf{target} \\
\midrule
coverage & $\sim$5$\times$ & full depth \\
window   & 25{,}000 bp & 25{,}000 bp \\
extra    & reference sequence & --- \\
\bottomrule
\end{tabular}
\end{center}

The reconstruction is judged against the deep track that was measured on the same
sample, so the ground truth is real data and not a simulation.

\section{Input}

The pipeline begins from \textbf{fragment centre counts}, one file per chromosome:

\begin{center}
\texttt{chr8\quad 4103877\quad 3\quad 2.8714}
\end{center}

Four tab-separated columns: chromosome, 1-based position, raw count, GC-corrected
count. Only covered positions are written; an absent position means zero coverage.
The pipeline reads the GC-corrected column.

These files are produced from BAMs by a separate project
(\texttt{cfdna\_tfp\_profiling}). That boundary is deliberate: this repository
starts from counts on disk and never reads a BAM.

The undersampled arm is made by randomly subsampling the BAM by read name, so both
ends of a fragment are kept or dropped together --- a half fragment has no midpoint.
The seed and the fraction are recorded, since they are the only things that make
that subsample reproducible.

Two further inputs describe the genome rather than the sample: the reference
sequence, and a \textbf{mappability list} of regions where short reads can be
placed reliably.

\section{Preprocessing}

\subsection{The smoothing chain}

Raw midpoint counts are spiky. Four steps in fixed order turn them into a
nucleosome track:

\begin{enumerate}
\item \textbf{Blacklist zeroing.} Positions in unreliable regions are set to zero.
\item \textbf{Whittaker--Eilers smoothing} ($\lambda = 1000$, order 2). Fits a
      smooth curve through the counts.
\item \textbf{Gaussian smoothing} ($\sigma = 30$ bp). Blurs at roughly the scale of
      a nucleosome.
\item \textbf{Subtract a running median} (window 375 bp). Removes slow drift in
      overall depth, leaving the nucleosome oscillation centred on zero.
\end{enumerate}

Both coverage levels go through the identical chain. Parts of it are known to be
redundant, and it is kept unchanged so that results stay comparable to earlier work.

\subsection{Smoothing happens when the data is read, not on disk}

The stored file holds \emph{raw} counts. Smoothing is applied by the dataloader
each time a window is read.

This costs training time --- it is the single largest cost per window --- and buys
two things: smoothing is destructive and some of its parameters are still open
questions, and one file can serve consumers that want the data differently.

\subsection{Windows and the halo}

The genome is cut into non-overlapping 25{,}000 bp windows, tiled from position~0.
Incomplete windows at a chromosome's end are dropped, because the model takes a
fixed input length.

Smoothing a window in isolation would corrupt its edges, since the filters need
neighbouring data that a bare window does not have. So each window is read with a
\textbf{400 bp halo} on each side, smoothed, and then cropped back to 25{,}000.
The halo comes from the stored array, never from padding: padding would reintroduce
the very edge artefact the halo exists to remove.

Gaussian smoothing reaches $4\sigma = 120$ bp and the running median reaches
187~bp, both comfortably inside 400. Whittaker--Eilers has no finite reach, so its
influence never becomes exactly zero --- but the residual it leaves is far below the
precision at which the counts are stored, and is therefore invisible in practice.

\subsection{Validity}

A position is \emph{valid} if it is mappable and no smoothing kernel that produced
it touched an unmappable base. In practice the mappability mask is grown by 120~bp
before use, matching the Gaussian's reach. Positions within 400~bp of a chromosome
end have no real halo and are excluded rather than repaired.

Validity is recorded per position and per window. It is currently used for
\emph{reporting}, not in the loss.

\subsection{Splitting}

Chromosomes are assigned whole to one split, so no window can appear on both sides:

\begin{center}
\begin{tabular}{ll}
\toprule
train & chr1--6, chr10--22 \\
val   & chr7 \\
test  & chr8, chr9 \\
\bottomrule
\end{tabular}
\end{center}

\section{The model}

A 1D U-Net: compress the signal three times, then expand it back, carrying detail
across at each level.

\subsection{Shape}

\begin{center}
\begin{tabular}{lll}
\toprule
\textbf{stage} & \textbf{length} & \textbf{channels} \\
\midrule
input        & $L$   & 6 \\
encoder 1    & $L/2$ & 64 \\
encoder 2    & $L/4$ & 128 \\
encoder 3    & $L/8$ & 256 \\
bottleneck   & $L/8$ & 256 \\
decoder 3--1 & back to $L$ & 256 $\to$ 128 $\to$ 64 \\
output       & $L$   & 1 \\
\bottomrule
\end{tabular}
\end{center}

Each stage halves the length and doubles the channels, so capacity stays roughly
constant while the view widens. Skip connections pass each encoder output to the
matching decoder, so fine detail is not lost in the compression. The output layer
is a $1\times1$ convolution with no activation: the target is centred near zero and
can be negative.

\subsection{Inside a block}

Every stage contains a residual block of four convolutions, kernel width 7, with
dilations 1, 2, 4 and 8. Dilation lets a block see 6, 12, 24 and 48~bp without
extra parameters --- local shape, transcription-factor footprint, and
sub-nucleosomal scale.

Two smaller choices matter. \textbf{Reflect padding} rather than zeros, because a
coverage window has no natural boundary and zero padding would look like an abrupt
loss of coverage. \textbf{GroupNorm} rather than BatchNorm, because the batch is
small when each example is 25{,}000 long.

The purely convolutional receptive field is 2{,}073~bp --- about ten nucleosomes.

\subsection{Three variants}

\begin{center}
\begin{tabular}{llr}
\toprule
\textbf{variant} & \textbf{description} & \textbf{parameters} \\
\midrule
V1 & coverage and sequence in one shared first convolution & 8{,}517{,}953 \\
V2 & coverage only; the sequence ablation                  & 8{,}515{,}393 \\
V3 & separate stream per input, joined before the encoder  & 8{,}545{,}665 \\
\bottomrule
\end{tabular}
\end{center}

V3 is the variant in current use. Coverage passes through its own 1D convolution;
sequence passes through a 2D convolution spanning all five bases at once and 11~bp
along the genome --- one full turn of the DNA helix, the periodicity at which
sequence influences nucleosome positioning. Each stream produces 32 channels, and
the two are concatenated to 64 before the encoder.

The point of the split is that one shared convolution has to serve a continuous
count and a categorical one-hot code at the same time; two streams let each be
read in its own terms.

\section{Loss and training}

The loss is \textbf{mean squared error} between the reconstruction and the deep
track, and nothing else.

Pearson correlation is computed on validation but carries no weight in the
objective. It is reported because MSE and correlation can move apart: a model can
lower MSE by shrinking its output toward zero without improving the shape at all.
Watching both makes that visible.

\begin{center}
\begin{tabular}{ll}
\toprule
optimiser & AdamW, learning rate $10^{-3}$, weight decay $10^{-2}$ \\
schedule  & halve the rate after 7 epochs without improvement, floor $10^{-5}$ \\
clipping  & gradient norm 1.0 \\
batch     & 16 windows \\
epochs    & 40 \\
seed      & 42 \\
\bottomrule
\end{tabular}
\end{center}

The gradient norm is logged \emph{before} clipping. If it sits near 1.0 the clip is
firing constantly and acting as a hidden reduction in learning rate rather than as
a safety net.

\textbf{The validity mask is available to the loss but switched off by default.}
Using it is a change to the objective, and mixing that with a change to the data
would leave neither attributable.

\section{Evaluation}

Evaluation runs on the held-out chromosomes, one window at a time.

\subsection{The comparison that matters}

Every metric is reported twice: once for the reconstruction, and once for the
\textbf{raw undersampled input} against the same target.

This is the point of the evaluation. A reconstruction correlating at 0.57 with the
truth means nothing on its own --- what matters is what the input already achieved
without any model. The gap between the two is the result.

\subsection{Metrics}

\textbf{Pearson correlation} measures agreement in shape and ignores amplitude.

\textbf{Mean squared error} measures agreement including amplitude.

\textbf{Peak distance} asks the biological question directly: are the nucleosomes
in the right places? Peaks are called with a frozen peak caller, and each peak is
matched to the nearest peak in the other track.

That last measurement has a trap. Matching each \emph{true} peak to the nearest
predicted one rewards predicting more peaks --- a track with a peak every few bases
scores perfectly. So the distance is measured in both directions:

\begin{center}
\begin{tabular}{ll}
\toprule
forward & one distance per true peak; asks \emph{were the real peaks found?} \\
reverse & one distance per predicted peak; asks \emph{are the predictions real?} \\
\bottomrule
\end{tabular}
\end{center}

Peak counts for all three tracks are reported alongside. An improvement should only
be quoted if it holds in both directions.

\textbf{A Wilcoxon signed-rank test} compares the two sets of peak distances window
by window. Paired, because each window supplies both arms.

\section{Output}

Training writes learning curves, the three best checkpoints by validation loss, and
the exact configuration that produced them. Each checkpoint carries its own
hyperparameters, so it can be reloaded without external notes.

Evaluation writes:

\begin{center}
\begin{tabular}{ll}
\toprule
\texttt{test\_summary.txt}        & the headline numbers \\
\texttt{test\_signal\_metrics.csv} & per window: MSE, both correlations, valid fraction \\
\texttt{test\_peak\_distances.csv} & per window: peak counts and both distances \\
\texttt{predictions.tsv.gz}       & the reconstructed track, as bedGraph \\
\texttt{ground\_truth.tsv.gz}     & the deep track, same format \\
\bottomrule
\end{tabular}
\end{center}

The per-window files are what allow the result to be examined rather than merely
quoted --- for example, split by mappability.

\section{Current results}

Held-out chromosomes 8 and 9, sample BH01 (healthy), 5$\times$ input:

\begin{center}
\begin{tabular}{lrr}
\toprule
 & \textbf{reconstruction} & \textbf{raw input} \\
\midrule
Pearson $r$                 & 0.5729 & 0.4234 \\
peak distance, forward      & 22 bp  & 34 bp  \\
peak distance, reverse      & 23 bp  & 31 bp  \\
peaks called (median)       & 126    & 115    \\
\bottomrule
\end{tabular}
\end{center}

The true track calls 122 peaks. Windows scored: 9{,}723. The reconstruction beats
its input in 98.8\% of windows, and wins in \emph{both} peak directions, so the
improvement is not an artefact of calling more peaks.

Performance depends on how mappable the region is:

\begin{center}
\begin{tabular}{lrrrr}
\toprule
\textbf{unmappable fraction} & \textbf{windows} & \textbf{recon $r$} & \textbf{input $r$} & \textbf{gain} \\
\midrule
$<$1\%      & 3{,}110 & 0.590 & 0.410 & $+0.179$ \\
1--10\%     & 4{,}997 & 0.584 & 0.419 & $+0.166$ \\
10--50\%    & 1{,}375 & 0.563 & 0.442 & $+0.126$ \\
$>$50\%     &    241  & 0.262 & 0.651 & $-0.358$ \\
\bottomrule
\end{tabular}
\end{center}

The last row is worth reading carefully. Where more than half the window is
unmappable, the \emph{input} appears to correlate better than anywhere else. That is
not signal: both coverage levels contain the same read-pileup artefacts, so they
agree with each other. The model does not reproduce those artefacts and scores
poorly there --- which is evidence that what it learned is nucleosome structure
rather than mapping error.

\section{Known limitations}

\begin{itemize}
\item The validity mask is not used in the loss. It is available and switched off.
\item Windows are kept if any position is valid, so a few almost entirely
      unmappable windows reach the test set.
\item An earlier, inherited data format used different window boundaries and a
      different position convention. Results from the two cannot be compared window
      by window; they agree at the level of summary metrics.
\item Reconstructed amplitude is systematically smaller than the truth. This is the
      expected behaviour of a squared-error objective under uncertainty, but it
      means amplitudes are not comparable between samples of different depth.
\end{itemize}

\end{document}
